\documentclass[11pt, a4paper]{article}
\usepackage[utf8]{inputenc}
\usepackage[margin=0.75in]{geometry}
\usepackage{amsmath, amssymb}
\usepackage{enumitem}
\usepackage{multicol}

\title{\textbf{NEET Biology: Glycolysis (EMP Pathway) — High-Yield Difficult Questions}}
\author{Target: NEET UG}
\date{}

\begin{document}

\maketitle

\section*{Questions}

\begin{enumerate}[label=\textbf{\arabic*.}]

\item Identify the incorrect statement regarding the universal metabolic pathway discovered by Gustav Embden, Otto Meyerhof, and J. Parnas:
\begin{enumerate}[label=(\Alph*)]
    \item It occurs in the cytoplasm of all living organisms, both obligate anaerobes and facultative aerobes.
    \item It neither requires molecular oxygen ($\text{O}_2$) nor releases carbon dioxide ($\text{CO}_2$).
    \item During the preparatory phase, two molecules of $\text{ATP}$ are consumed to form a symmetric hexose bisphosphate.
    \item Every enzymatic step in the pathway is catalyzed by a distinct isoform of dehydrogenase and kinase.
\end{enumerate}

\item Which of the following pairs of enzymes in the EMP pathway are allosterically regulated and catalyze irreversible reactions under intracellular physiological conditions?
\begin{enumerate}[label=(\Alph*)]
    \item Phosphohexose isomerase, Aldolase, Enolase
    \item Hexokinase, Phosphofructokinase-1, Pyruvate kinase
    \item Triose phosphate isomerase, Phosphoglycerate kinase, Mutase
    \item Glucokinase, Glyceraldehyde-3-phosphate dehydrogenase, Enolase
\end{enumerate}

\item Consider the conversion of $1\text{ mole}$ of sucrose into pyruvic acid via glycolysis. What is the net gain of $\text{ATP}$ (via substrate-level phosphorylation) and $\text{NADH} + \text{H}^+$, respectively?
\begin{enumerate}[label=(\Alph*)]
    \item $2\text{ ATP}$ and $2\text{ NADH} + \text{H}^+$
    \item $4\text{ ATP}$ and $4\text{ NADH} + \text{H}^+$
    \item $8\text{ ATP}$ and $4\text{ NADH} + \text{H}^+$
    \item $4\text{ ATP}$ and $2\text{ NADH} + \text{H}^+$
\end{enumerate}

\item \textbf{Assertion (A):} The conversion of 3-phosphoglyceraldehyde ($\text{PGAL}$) to 1,3-bisphosphoglycerate ($\text{BPGA}$) is coupled with the reduction of $\text{NAD}^+$ to $\text{NADH} + \text{H}^+$. \\
\textbf{Reason (R):} The inorganic phosphate ($\text{H}_3\text{PO}_4$) added at the $\text{C}_1$ position of $\text{PGAL}$ is donated directly by an $\text{ATP}$ molecule.
\begin{enumerate}[label=(\Alph*)]
    \item Both (A) and (R) are true and (R) is the correct explanation of (A).
    \item Both (A) and (R) are true but (R) is not the correct explanation of (A).
    \item (A) is true but (R) is false.
    \item Both (A) and (R) are false.
\end{enumerate}

\item If a glucose molecule is radiolabeled with $^{14}\text{C}$ exclusively at Carbon-1 ($\text{C}_1$), which carbon(s) of the two resulting pyruvate molecules will carry the radioactive label?
\begin{enumerate}[label=(\Alph*)]
    \item The carboxyl carbon ($\text{C}_1$) of one pyruvate molecule only.
    \item The methyl carbon ($\text{C}_3$) of one pyruvate molecule only.
    \item The carbonyl carbon ($\text{C}_2$) of both pyruvate molecules.
    \item The methyl carbon ($\text{C}_3$) of both pyruvate molecules.
\end{enumerate}

\item Match Column-I (Glycolytic Enzyme) with Column-II (Type of Reaction / Cofactor Requirement):
\begin{center}
\begin{tabular}{|l|l|}
\hline
\textbf{Column-I} & \textbf{Column-II} \\
\hline
P. Aldolase & 1. Dehydration; requires $\text{Mg}^{2+}$ \\
Q. Enolase & 2. Reversible cleavage of a 6-C ketose into two 3-C trioses \\
R. Phosphoglycerate mutase & 3. Intramolecular shift of phosphoryl group \\
S. Glyceraldehyde-3-P dehydrogenase & 4. Oxidative phosphorylation at substrate level \\
\hline
\end{tabular}
\end{center}
\begin{enumerate}[label=(\Alph*)]
    \item P-2, Q-1, R-3, S-4
    \item P-2, Q-3, R-1, S-4
    \item P-4, Q-1, R-3, S-2
    \item P-1, Q-2, R-4, S-3
\end{enumerate}

\item In plants, sucrose is converted into glucose and fructose by the enzyme \_\_\_\_\_\_\_\_, and these two monosaccharides enter the glycolytic pathway after phosphorylation by \_\_\_\_\_\_\_\_.
\begin{enumerate}[label=(\Alph*)]
    \item Sucrase; Phosphofructokinase
    \item Invertase; Hexokinase
    \item Zymase; Glucokinase
    \item Amylase; Aldolase
\end{enumerate}

\item During the payoff phase of glycolysis, water ($\text{H}_2\text{O}$) is eliminated during which specific enzymatic transformation?
\begin{enumerate}[label=(\Alph*)]
    \item 1,3-bisphosphoglycerate $\rightarrow$ 3-phosphoglycerate
    \item 3-phosphoglycerate $\rightarrow$ 2-phosphoglycerate
    \item 2-phosphoglycerate $\rightarrow$ Phosphoenolpyruvate
    \item Phosphoenolpyruvate $\rightarrow$ Pyruvic acid
\end{enumerate}

\item Which of the following inhibitors specifically blocks glycolysis by competing with inorganic phosphate ($\text{P}_i$) during the oxidation of glyceraldehyde-3-phosphate, thereby preventing net $\text{ATP}$ synthesis without stopping the oxidation step itself?
\begin{enumerate}[label=(\Alph*)]
    \item Fluoride ($\text{F}^-$)
    \item Iodoacetate
    \item Arsenate ($\text{AsO}_4^{3-}$)
    \item Malonate
\end{enumerate}

\item Sodium fluoride ($\text{NaF}$) is routinely added to blood collection vials for blood glucose estimation because fluoride ions inhibit glycolysis by chelating $\text{Mg}^{2+}$ required by the enzyme:
\begin{enumerate}[label=(\Alph*)]
    \item Hexokinase
    \item Enolase
    \item Pyruvate kinase
    \item Phosphoglycerate kinase
\end{enumerate}

\item How many total $\text{ATP}$ equivalents are generated when $1\text{ molecule}$ of glucose is completely oxidized to pyruvic acid under aerobic conditions, assuming the malate-aspartate shuttle operates to transport cytosolic reducing equivalents into the mitochondrial matrix?
\begin{enumerate}[label=(\Alph*)]
    \item $2\text{ ATP}$
    \item $6\text{ ATP}$
    \item $8\text{ ATP}$
    \item $38\text{ ATP}$
\end{enumerate}

\item Identify the correct sequence of intermediates formed during the 10-step glycolytic pathway:
\begin{enumerate}[label=(\Alph*)]
    \item $\text{F-6-P} \rightarrow \text{F-1,6-BP} \rightarrow \text{DHAP} \leftrightarrow \text{PGAL} \rightarrow \text{1,3-BPGA} \rightarrow \text{3-PGA} \rightarrow \text{2-PGA} \rightarrow \text{PEP}$
    \item $\text{F-6-P} \rightarrow \text{F-1,6-BP} \rightarrow \text{PGAL} \rightarrow \text{3-PGA} \rightarrow \text{1,3-BPGA} \rightarrow \text{2-PGA} \rightarrow \text{PEP}$
    \item $\text{G-6-P} \rightarrow \text{F-1,6-BP} \rightarrow \text{F-6-P} \rightarrow \text{DHAP} \rightarrow \text{1,3-BPGA} \rightarrow \text{2-PGA} \rightarrow \text{PEP}$
    \item $\text{G-6-P} \rightarrow \text{F-6-P} \rightarrow \text{F-1,6-BP} \rightarrow \text{PGAL} \rightarrow \text{2-PGA} \rightarrow \text{3-PGA} \rightarrow \text{PEP}$
\end{enumerate}

\item Which step in glycolysis represents the ``committed step'' (rate-limiting pacemaker step) of the pathway?
\begin{enumerate}[label=(\Alph*)]
    \item Phosphorylation of glucose to glucose-6-phosphate
    \item Isomerization of glucose-6-phosphate to fructose-6-phosphate
    \item Phosphorylation of fructose-6-phosphate to fructose-1,6-bisphosphate
    \item Conversion of phosphoenolpyruvate to pyruvate
\end{enumerate}

\item \textbf{Statement I:} In glycolysis, substrate-level phosphorylation occurs twice per triose phosphate molecule metabolized. \\
\textbf{Statement II:} Both substrate-level phosphorylation steps in glycolysis are catalyzed by reversible kinase enzymes.
\begin{enumerate}[label=(\Alph*)]
    \item Both Statement I and Statement II are correct.
    \item Both Statement I and Statement II are incorrect.
    \item Statement I is correct, but Statement II is incorrect.
    \item Statement I is incorrect, but Statement II is correct.
\end{enumerate}

\item What is the ratio of total $\text{ATP}$ consumed to total $\text{ATP}$ synthesized (gross) during the glycolytic breakdown of one molecule of glucose-6-phosphate derived directly from glycogenolysis?
\begin{enumerate}[label=(\Alph*)]
    \item $1 : 2$
    \item $1 : 4$
    \item $2 : 4$
    \item $1 : 3$
\end{enumerate}

\item If a glucose molecule is labeled with $^{14}\text{C}$ at Carbon-3 ($\text{C}_3$) and Carbon-4 ($\text{C}_4$), and subjected to alcoholic fermentation via glycolysis, where will the $^{14}\text{C}$ label appear at the end of the process?
\begin{enumerate}[label=(\Alph*)]
    \item Exclusively in ethanol
    \item Exclusively in $\text{CO}_2$
    \item Equally distributed between ethanol and $\text{CO}_2$
    \item Exclusively in acetaldehyde
\end{enumerate}

\item Which of the following glycolytic intermediates is an allosteric feed-forward activator of pyruvate kinase in mammalian liver cells?
\begin{enumerate}[label=(\Alph*)]
    \item Glucose-6-phosphate
    \item Fructose-1,6-bisphosphate
    \item 3-Phosphoglycerate
    \item Phosphoenolpyruvate
\end{enumerate}

\item Why does the isomerization of glucose-6-phosphate (an aldose) to fructose-6-phosphate (a ketose) occur prior to the second phosphorylation step?
\begin{enumerate}[label=(\Alph*)]
    \item To prevent glucose-6-phosphate from diffusing out of the plasma membrane.
    \item To expose the primary hydroxyl group at $\text{C}_1$ for phosphorylation and facilitate subsequent symmetric cleavage between $\text{C}_3$ and $\text{C}_4$.
    \item To generate an immediate high-energy thioester bond.
    \item Because hexokinase cannot bind to pyranose rings.
\end{enumerate}

\item How many total phosphate groups (both inorganic and organic) are present in the products of the reaction catalyzed by Glyceraldehyde-3-phosphate dehydrogenase per glucose molecule?
\begin{enumerate}[label=(\Alph*)]
    \item 2
    \item 3
    \item 4
    \item 6
\end{enumerate}

\item Which of the following statements accurately distinguishes between Hexokinase and Glucokinase?
\begin{enumerate}[label=(\Alph*)]
    \item Hexokinase has a high $K_m$ (low affinity) and is inhibited by glucose-6-phosphate, whereas glucokinase has a low $K_m$ (high affinity).
    \item Hexokinase has a low $K_m$ (high affinity) and is allosterically inhibited by glucose-6-phosphate, whereas glucokinase is liver-specific, has a high $K_m$, and is not inhibited by glucose-6-phosphate.
    \item Both enzymes are found exclusively in skeletal muscles and require $\text{Mn}^{2+}$ instead of $\text{Mg}^{2+}$.
    \item Glucokinase catalyzes the phosphorylation of fructose and mannose more efficiently than glucose.
\end{enumerate}

\item In erythrocytes (mature red blood cells), 1,3-bisphosphoglycerate can be diverted via the Rapoport-Luebering shunt to form 2,3-bisphosphoglycerate ($2,3\text{-BPG}$). What is the energetic consequence of this bypass per molecule of $\text{BPG}$ diverted?
\begin{enumerate}[label=(\Alph*)]
    \item Net gain of $1\text{ extra ATP}$
    \item Loss of $1\text{ ATP}$ synthesis step (catalyzed by phosphoglycerate kinase)
    \item Loss of $2\text{ ATP}$ synthesis steps
    \item Increased production of $\text{NADH} + \text{H}^+$
\end{enumerate}

\item Which of the following glycolytic intermediates possesses the highest standard free energy of hydrolysis ($\Delta G^{\circ\prime}$) of its phosphate group, exceeding even that of $\text{ATP}$?
\begin{enumerate}[label=(\Alph*)]
    \item Glucose-6-phosphate
    \item Fructose-1,6-bisphosphate
    \item 3-Phosphoglycerate
    \item Phosphoenolpyruvate ($\text{PEP}$)
\end{enumerate}

\item Consider the following statements regarding Triose Phosphate Isomerase (TPI):
\begin{enumerate}[label=(\roman*)]
    \item It interconverts dihydroxyacetone phosphate ($\text{DHAP}$) and glyceraldehyde-3-phosphate ($\text{PGAL}$).
    \item At equilibrium, $\text{PGAL}$ accounts for over $95\%$ of the triose phosphate pool.
    \item Without TPI, the net $\text{ATP}$ yield of glycolysis from one glucose molecule would drop to zero.
    \item It is considered a ``catalytically perfect'' enzyme whose rate is limited only by substrate diffusion.
\end{enumerate}
Which of the above statements are correct?
\begin{enumerate}[label=(\Alph*)]
    \item (i), (ii), and (iii)
    \item (i), (iii), and (iv)
    \item (ii) and (iv) only
    \item (i) and (iii) only
\end{enumerate}

\item During glycolysis, how many molecules of $\text{H}^+$ (protons) are released into the cytosol during the oxidation of two molecules of glyceraldehyde-3-phosphate to two molecules of 1,3-bisphosphoglycerate?
\begin{enumerate}[label=(\Alph*)]
    \item 1
    \item 2
    \item 4
    \item 0
\end{enumerate}

\item What is the Pasteur effect in cellular respiration?
\begin{enumerate}[label=(\Alph*)]
    \item Stimulation of glycolysis by high concentrations of oxygen.
    \item Inhibition of glycolysis and decreased glucose consumption when anaerobic cells are exposed to oxygen.
    \item Conversion of pyruvate to ethanol in the presence of excess $\text{ATP}$.
    \item Inhibition of the Krebs cycle by high levels of fructose-1,6-bisphosphate.
\end{enumerate}

\item Identify the two glycolytic reactions that involve the cleavage or formation of carbon-carbon ($\text{C}-\text{C}$) bonds:
\begin{enumerate}[label=(\Alph*)]
    \item Hexokinase and Pyruvate kinase
    \item Aldolase only (no $\text{C}-\text{C}$ bond is formed in glycolysis)
    \item Phosphofructokinase-1 and Enolase
    \item Aldolase and Glyceraldehyde-3-phosphate dehydrogenase
\end{enumerate}

\item Which of the following divalent cations is an essential cofactor for almost all kinase and enolase reactions in the EMP pathway?
\begin{enumerate}[label=(\Alph*)]
    \item $\text{Fe}^{2+}$
    \item $\text{Ca}^{2+}$
    \item $\text{Mg}^{2+}$
    \item $\text{Zn}^{2+}$
\end{enumerate}

\item If a cell is treated with a drug that completely inhibits Lactate Dehydrogenase ($\text{LDH}$) under strictly anaerobic conditions, glycolysis will quickly halt because:
\begin{enumerate}[label=(\Alph*)]
    \item Pyruvate accumulates and directly inhibits hexokinase.
    \item Cytosolic $\text{NAD}^+$ is depleted, halting the glyceraldehyde-3-phosphate dehydrogenase reaction.
    \item $\text{ATP}$ levels rise excessively, inhibiting phosphofructokinase-1.
    \item Phosphoenolpyruvate cannot be converted into pyruvate.
\end{enumerate}

\item \textbf{Assertion (A):} In glycolysis, the first phosphorylation of glucose traps it inside the cell. \\
\textbf{Reason (R):} Plasma membranes lack specific transporters for phosphorylated sugars like glucose-6-phosphate, and the negative charge on the phosphate group prevents passive diffusion across the lipid bilayer.
\begin{enumerate}[label=(\Alph*)]
    \item Both (A) and (R) are true and (R) is the correct explanation of (A).
    \item Both (A) and (R) are true but (R) is not the correct explanation of (A).
    \item (A) is true but (R) is false.
    \item Both (A) and (R) are false.
\end{enumerate}

\item How many redox reactions occur during the entire 10-step glycolytic conversion of glucose to pyruvate?
\begin{enumerate}[label=(\Alph*)]
    \item Zero
    \item One (occurring twice per glucose molecule)
    \item Two (occurring twice per glucose molecule)
    \item Three
\end{enumerate}

\item Which intermediate of glycolysis serves as the direct precursor for the synthesis of the glycerol backbone of triacylglycerols (fats) in adipose tissue and liver?
\begin{enumerate}[label=(\Alph*)]
    \item 3-Phosphoglycerate
    \item Dihydroxyacetone phosphate ($\text{DHAP}$)
    \item Phosphoenolpyruvate
    \item Glucose-6-phosphate
\end{enumerate}

\item In the reaction catalyzed by Phosphoglycerate mutase, 3-phosphoglycerate is converted into 2-phosphoglycerate via which transient enzyme-bound intermediate?
\begin{enumerate}[label=(\Alph*)]
    \item 1,3-bisphosphoglycerate
    \item 2,3-bisphosphoglycerate
    \item 1,2-bisphosphoglycerate
    \item Phosphoenolpyruvate
\end{enumerate}

\item Match the metabolic fate of pyruvate with the organism/condition in which it predominantly occurs:
\begin{center}
\begin{tabular}{|l|l|}
\hline
\textbf{Column-I (Condition/Organism)} & \textbf{Column-II (End Product \& Key Enzyme)} \\
\hline
P. Vigorous exercise in human skeletal muscle & 1. Ethanol + $\text{CO}_2$; Pyruvate decarboxylase \& Alcohol dehydrogenase \\
Q. *Saccharomyces cerevisiae* under hypoxia & 2. Acetyl-CoA + $\text{CO}_2$; Pyruvate dehydrogenase complex \\
R. Aerobic respiration in eukaryotic mitochondria & 3. Lactic acid (no $\text{CO}_2$); Lactate dehydrogenase \\
\hline
\end{tabular}
\end{center}
\begin{enumerate}[label=(\Alph*)]
    \item P-3, Q-1, R-2
    \item P-1, Q-3, R-2
    \item P-3, Q-2, R-1
    \item P-2, Q-1, R-3
\end{enumerate}

\item Which of the following enzymes in glycolysis belongs to the class \textbf{Lyases} according to the IUBMB enzyme classification?
\begin{enumerate}[label=(\Alph*)]
    \item Hexokinase and Pyruvate kinase
    \item Phosphohexose isomerase and Triose phosphate isomerase
    \item Aldolase and Enolase
    \item Phosphoglycerate kinase and Phosphoglycerate mutase
\end{enumerate}

\item High cellular concentrations of which molecule act as a potent allosteric \textbf{inhibitor} of Phosphofructokinase-1 ($\text{PFK-1}$), signaling energy abundance from the Krebs cycle?
\begin{enumerate}[label=(\Alph*)]
    \item $\text{AMP}$
    \item $\text{ADP}$
    \item Citrate
    \item Fructose-2,6-bisphosphate
\end{enumerate}

\item Calculate the total number of carbon atoms in all the intermediate molecules formed simultaneously right after the aldolase cleavage of one molecule of fructose-1,6-bisphosphate:
\begin{enumerate}[label=(\Alph*)]
    \item 3 carbons
    \item 6 carbons (two 3-carbon molecules)
    \item 12 carbons
    \item 4 carbons
\end{enumerate}

\item What is the primary reason why the conversion of phosphoenolpyruvate ($\text{PEP}$) to pyruvate is thermodynamically highly exergonic ($\Delta G^{\circ\prime} = -31.4\text{ kJ/mol}$)?
\begin{enumerate}[label=(\Alph*)]
    \item Pyruvate immediately enters the mitochondrion, pulling the reaction forward.
    \item The unstable enol form of pyruvate rapidly tautomerizes into the much more stable keto form.
    \item $\text{PEP}$ contains two high-energy anhydride bonds.
    \item Pyruvate kinase consumes $\text{NADH}$ during the transfer.
\end{enumerate}

\item In the payoff phase of glycolysis, how many molecules of water are produced per molecule of glucose metabolized to pyruvate?
\begin{enumerate}[label=(\Alph*)]
    \item 1
    \item 2
    \item 4
    \item 0
\end{enumerate}

\item Which of the following statements about glycolysis is \textbf{false}?
\begin{enumerate}[label=(\Alph*)]
    \item Partial oxidation of glucose occurs without the help of oxygen.
    \item In anaerobic organisms, glycolysis is the only process in respiration.
    \item Pyruvic acid is the key product of glycolysis, and its metabolic fate depends on cellular need.
    \item Glycolysis occurs in the mitochondrial matrix in eukaryotes and cytosol in prokaryotes.
\end{enumerate}

\item During the conversion of 1,3-bisphosphoglycerate to 3-phosphoglycerate, which phosphate group is transferred to $\text{ADP}$ to form $\text{ATP}$?
\begin{enumerate}[label=(\Alph*)]
    \item The low-energy phosphate ester at $\text{C}_3$
    \item The high-energy acyl phosphate at $\text{C}_1$
    \item Both phosphates are transferred sequentially
    \item An inorganic phosphate from the cytosol
\end{enumerate}

\item If $5\text{ molecules}$ of glucose undergo glycolysis, how many total $\text{ATP}$ molecules are \textbf{synthesized} (gross yield) and how many are \textbf{consumed}, respectively?
\begin{enumerate}[label=(\Alph*)]
    \item $10\text{ synthesized}$, $5\text{ consumed}$
    \item $20\text{ synthesized}$, $10\text{ consumed}$
    \item $10\text{ synthesized}$, $10\text{ consumed}$
    \item $40\text{ synthesized}$, $10\text{ consumed}$
\end{enumerate}

\item Which amino acid is synthesized directly via transamination of pyruvate, linking glycolysis to amino acid metabolism?
\begin{enumerate}[label=(\Alph*)]
    \item Serine
    \item Alanine
    \item Glycine
    \item Aspartate
\end{enumerate}

\item What is the oxidation state change of the aldehyde carbon ($\text{C}_1$) of glyceraldehyde-3-phosphate when it is converted into 1,3-bisphosphoglycerate?
\begin{enumerate}[label=(\Alph*)]
    \item Oxidized from an aldehyde to a carboxylic acid derivative (acyl phosphate)
    \item Reduced from an aldehyde to a primary alcohol
    \item Oxidized from a ketone to an ester
    \item No change in oxidation state occurs
\end{enumerate}

\item \textbf{Statement I:} Fructose-2,6-bisphosphate is an intermediate of the glycolytic pathway. \\
\textbf{Statement II:} Fructose-2,6-bisphosphate is the most potent allosteric activator of phosphofructokinase-1 ($\text{PFK-1}$) in hepatocytes.
\begin{enumerate}[label=(\Alph*)]
    \item Both Statement I and Statement II are correct.
    \item Both Statement I and Statement II are incorrect.
    \item Statement I is correct, but Statement II is incorrect.
    \item Statement I is incorrect, but Statement II is correct.
\end{enumerate}

\item In a hypothetical mutant cell where Enolase is completely non-functional, which glycolytic intermediate would accumulate to the highest concentration if glucose is supplied continuously?
\begin{enumerate}[label=(\Alph*)]
    \item Phosphoenolpyruvate
    \item 2-Phosphoglycerate
    \item Pyruvic acid
    \item Fructose-1,6-bisphosphate
\end{enumerate}

\end{enumerate}

\newpage

\section*{Answer Key}

\begin{multicols}{5}
\begin{enumerate}[label=\textbf{\arabic*.}]
    \item (D)
    \item (B)
    \item (B)
    \item (C)
    \item (B)
    \item (A)
    \item (B)
    \item (C)
    \item (C)
    \item (B)
    \item (C)
    \item (A)
    \item (C)
    \item (C)
    \item (B)
    \item (B)
    \item (B)
    \item (B)
    \item (C)
    \item (B)
    \item (B)
    \item (D)
    \item (B)
    \item (B)
    \item (B)
    \item (B)
    \item (C)
    \item (B)
    \item (A)
    \item (B)
    \item (B)
    \item (B)
    \item (A)
    \item (C)
    \item (C)
    \item (B)
    \item (B)
    \item (B)
    \item (D)
    \item (B)
    \item (B)
    \item (B)
    \item (A)
    \item (D)
    \item (B)
\end{enumerate}
\end{multicols}

\section*{Detailed Solutions}

\begin{enumerate}[label=\textbf{\arabic*.}]

\item \textbf{(D)} Only one step in glycolysis involves a dehydrogenase (Glyceraldehyde-3-phosphate dehydrogenase), and four steps involve kinases (Hexokinase, PFK-1, Phosphoglycerate kinase, Pyruvate kinase). The remaining steps are catalyzed by isomerases, an aldolase, a mutase, and an enolase.

\item \textbf{(B)} Hexokinase (Step 1), Phosphofructokinase-1 (Step 3), and Pyruvate kinase (Step 10) catalyze reactions with large negative $\Delta G$, making them physiologically irreversible and the primary sites of allosteric regulation.

\item \textbf{(B)} One mole of sucrose ($12\text{C}$) is hydrolyzed by invertase into $1\text{ mole}$ of glucose ($6\text{C}$) and $1\text{ mole}$ of fructose ($6\text{C}$). Each hexose yields a net gain of $2\text{ ATP}$ and $2\text{ NADH} + \text{H}^+$. Therefore, $1\text{ mole}$ of sucrose yields $4\text{ ATP}$ and $4\text{ NADH} + \text{H}^+$.

\item \textbf{(C)} Assertion (A) is true: $\text{PGAL}$ is oxidized while $\text{NAD}^+$ is reduced to $\text{NADH} + \text{H}^+$. Reason (R) is false because the phosphate added at $\text{C}_1$ comes from free inorganic phosphate ($\text{P}_i$ / $\text{H}_3\text{PO}_4$) in the cytosol, not from $\text{ATP}$.

\item \textbf{(B)} During aldolase cleavage, $\text{C}_1, \text{C}_2, \text{C}_3$ of fructose-1,6-bisphosphate form $\text{DHAP}$ ($\text{C}_1$ becomes the phosphorylated $\text{CH}_2\text{OPO}_3^{2-}$ group), while $\text{C}_4, \text{C}_5, \text{C}_6$ form $\text{PGAL}$. When $\text{DHAP}$ isomerizes to $\text{PGAL}$, its $\text{C}_1$ becomes $\text{C}_3$ of $\text{PGAL}$, which ultimately becomes the methyl carbon ($\text{C}_3$) of one pyruvate molecule.

\item \textbf{(A)} Aldolase splits $\text{F-1,6-BP}$ into $\text{DHAP}$ and $\text{PGAL}$ (P-2). Enolase removes $\text{H}_2\text{O}$ from $\text{2-PGA}$ using $\text{Mg}^{2+}$ (Q-1). Phosphoglycerate mutase shifts the phosphate from $\text{C}_3$ to $\text{C}_2$ (R-3). $\text{GAPDH}$ couples oxidation with inorganic phosphorylation (S-4).

\item \textbf{(B)} According to NCERT, sucrose is converted into glucose and fructose by the enzyme \textbf{invertase}, and these two monosaccharides readily enter the glycolytic pathway after phosphorylation by \textbf{hexokinase}.

\item \textbf{(C)} Enolase catalyzes the dehydration (removal of $\text{H}_2\text{O}$) of 2-phosphoglycerate to form phosphoenolpyruvate ($\text{PEP}$).

\item \textbf{(C)} Arsenate ($\text{AsO}_4^{3-}$) mimics inorganic phosphate ($\text{P}_i$). $\text{GAPDH}$ incorporates arsenate to form 1-arseno-3-phosphoglycerate, which spontaneously hydrolyzes to 3-phosphoglycerate, bypassing the $\text{ATP}$-generating step of phosphoglycerate kinase.

\item \textbf{(B)} Fluoride ions ($\text{F}^-$) react with $\text{Mg}^{2+}$ and phosphate to form a fluorophosphate complex that tightly binds and inhibits \textbf{enolase}, stopping glucose consumption in blood samples.

\item \textbf{(C)} Glycolysis yields a net of $2\text{ ATP}$ directly via substrate-level phosphorylation plus $2\text{ NADH} + \text{H}^+$. Via the malate-aspartate shuttle, each cytosolic $\text{NADH}$ yields $3\text{ ATP}$ in the electron transport chain ($2 \times 3 = 6\text{ ATP}$). Total = $2 + 6 = 8\text{ ATP}$.

\item \textbf{(A)} The exact sequence of intermediates after Glucose-6-phosphate is: Fructose-6-phosphate $\rightarrow$ Fructose-1,6-bisphosphate $\rightarrow$ $\text{DHAP} \leftrightarrow \text{PGAL} \rightarrow$ 1,3-bisphosphoglycerate $\rightarrow$ 3-phosphoglycerate $\rightarrow$ 2-phosphoglycerate $\rightarrow$ Phosphoenolpyruvate.

\item \textbf{(C)} Phosphofructokinase-1 ($\text{PFK-1}$) catalyzes the phosphorylation of fructose-6-phosphate to fructose-1,6-bisphosphate. Since glucose-6-phosphate can enter other pathways (glycogenesis, pentose phosphate pathway), $\text{PFK-1}$ is the first committed, rate-limiting pacemaker enzyme of glycolysis.

\item \textbf{(C)} Statement I is correct: per triose phosphate ($\text{PGAL}$), $\text{ATP}$ is formed once at step 7 (phosphoglycerate kinase) and once at step 10 (pyruvate kinase). Statement II is incorrect: phosphoglycerate kinase is reversible, whereas pyruvate kinase is physiologically irreversible.

\item \textbf{(B)} Glycogenolysis (via glycogen phosphorylase) produces glucose-1-phosphate without consuming $\text{ATP}$, which isomerizes to glucose-6-phosphate. Thus, only $1\text{ ATP}$ is consumed (at the $\text{PFK-1}$ step), while $4\text{ ATP}$ are synthesized in the payoff phase. Ratio = $1 : 4$.

\item \textbf{(B)} $\text{C}_3$ and $\text{C}_4$ of glucose become the carboxyl carbons ($\text{C}_1$) of the two pyruvate molecules. During alcoholic fermentation, pyruvate decarboxylase removes this carboxyl group as $\text{CO}_2$, leaving the remaining 2-carbon acetaldehyde (and subsequently ethanol) unlabeled.

\item \textbf{(B)} Fructose-1,6-bisphosphate acts as a feed-forward allosteric activator of pyruvate kinase, ensuring that the downstream payoff phase accelerates when upstream intermediates accumulate.

\item \textbf{(B)} Isomerizing the 6-membered aldose ring ($\text{G-6-P}$) into a 5-membered ketose ring ($\text{F-6-P}$) exposes the $-\text{CH}_2\text{OH}$ group at $\text{C}_1$ for phosphorylation and places the carbonyl group at $\text{C}_2$, enabling aldolase to cleave the $\text{C}_3-\text{C}_4$ bond cleanly into two 3-carbon trioses.

\item \textbf{(C)} Per glucose molecule, two molecules of $\text{PGAL}$ are converted into two molecules of 1,3-bisphosphoglycerate. Each 1,3-bisphosphoglycerate molecule contains $2\text{ phosphate groups}$, giving $2 \times 2 = 4\text{ phosphate groups}$ in total.

\item \textbf{(B)} Hexokinase (present in most tissues) has a low $K_m$ (high affinity for glucose) and is feedback-inhibited by $\text{G-6-P}$. Glucokinase (hexokinase IV, in liver and pancreatic $\beta$-cells) has a high $K_m$ (low affinity, active after meals) and is not inhibited by $\text{G-6-P}$.

\item \textbf{(B)} Diverting 1,3-bisphosphoglycerate to $2,3\text{-BPG}$ via bisphosphoglycerate mutase and subsequent hydrolysis to 3-phosphoglycerate bypasses the phosphoglycerate kinase step, forfeiting $1\text{ ATP}$ molecule per triose diverted (used in RBCs to regulate hemoglobin oxygen affinity).

\item \textbf{(D)} Phosphoenolpyruvate ($\text{PEP}$) has a $\Delta G^{\circ\prime}$ of hydrolysis of $-61.9\text{ kJ/mol}$ (compared to $-30.5\text{ kJ/mol}$ for $\text{ATP}$), making it the highest-energy phosphate compound in glycolysis.

\item \textbf{(B)} Statements (i), (iii), and (iv) are correct. Statement (ii) is incorrect because at equilibrium, $\text{DHAP}$ actually makes up $\sim 96\%$ of the pool, though rapid consumption of $\text{PGAL}$ by $\text{GAPDH}$ pulls the reaction forward. Without TPI, only $1\text{ PGAL}$ would enter the payoff phase ($2\text{ ATP}$ produced minus $2\text{ ATP}$ invested = $0\text{ net ATP}$).

\item \textbf{(B)} The balanced stoichiometry for one $\text{PGAL}$ oxidation is: $\text{PGAL} + \text{NAD}^+ + \text{P}_i \rightarrow \text{1,3-BPGA} + \text{NADH} + \text{H}^+$. For two molecules of $\text{PGAL}$ (from one glucose), $2\text{ H}^+$ ions are released.

\item \textbf{(B)} Discovered by Louis Pasteur in yeast, the Pasteur effect is the rapid drop in glucose consumption and inhibition of glycolysis (specifically via $\text{ATP}$/citrate inhibition of $\text{PFK-1}$) when anaerobic cells are shifted to aerobic conditions.

\item \textbf{(B)} Aldolase catalyzes the cleavage of the $\text{C}_3-\text{C}_4$ carbon-carbon bond of fructose-1,6-bisphosphate. No carbon-carbon bond is formed during glycolysis.

\item \textbf{(C)} $\text{Mg}^{2+}$ forms complexes with the negatively charged phosphoryl oxygens of $\text{ATP}$/$\text{ADP}$ and substrates, neutralizing charge repulsion and stabilizing transition states in kinases and enolase.

\item \textbf{(B)} Under anaerobic conditions, $\text{LDH}$ regenerates $\text{NAD}^+$ from $\text{NADH} + \text{H}^+$ by reducing pyruvate to lactate. Without $\text{LDH}$, all cytosolic $\text{NAD}^+$ is trapped as $\text{NADH}$, starving $\text{GAPDH}$ of its required electron acceptor.

\item \textbf{(A)} Phosphorylation adds bulky, negatively charged phosphate groups to glucose. Because cell membranes lack transport proteins for $\text{G-6-P}$ and charged molecules cannot cross the hydrophobic lipid bilayer, glucose is metabolically trapped inside the cell.

\item \textbf{(B)} Only Step 6 (catalyzed by Glyceraldehyde-3-phosphate dehydrogenase) is an oxidation-reduction reaction. Because two molecules of $\text{PGAL}$ are processed per glucose, this single type of redox reaction occurs twice per glucose.

\item \textbf{(B)} Dihydroxyacetone phosphate ($\text{DHAP}$) is reduced by glycerol-3-phosphate dehydrogenase (using $\text{NADH}$) to glycerol-3-phosphate, which forms the backbone of triglycerides and phospholipids.

\item \textbf{(B)} Phosphoglycerate mutase transfers a phosphate from an active-site histidine to the $\text{C}_2$ position of 3-phosphoglycerate, forming a transient enzyme-bound **2,3-bisphosphoglycerate** intermediate before removing the $\text{C}_3$ phosphate.

\item \textbf{(A)} Vigorous human muscle produces lactic acid via $\text{LDH}$ without $\text{CO}_2$ release (P-3). Yeast (*Saccharomyces cerevisiae*) performs alcoholic fermentation releasing ethanol and $\text{CO}_2$ (Q-1). Aerobic mitochondria decarboxylate pyruvate to Acetyl-CoA via the $\text{PDH}$ complex (R-2).

\item \textbf{(C)} Lyases cleave bonds without hydrolysis or oxidation, often forming double bonds or rings. Aldolase cleaves $\text{C}-\text{C}$ bonds, and Enolase removes $\text{H}_2\text{O}$ to form a double bond in $\text{PEP}$; both are Lyases.

\item \textbf{(C)} High citrate levels signal that the Krebs cycle is saturated with biosynthetic precursors and energy; citrate exits the mitochondrion and allosterically inhibits $\text{PFK-1}$ (enhancing the inhibitory effect of $\text{ATP}$).

\item \textbf{(B)} Aldolase splits one 6-carbon fructose-1,6-bisphosphate molecule into one 3-carbon $\text{DHAP}$ and one 3-carbon $\text{PGAL}$ molecule, totaling $3 + 3 = 6\text{ carbon atoms}$.

\item \textbf{(B)} Removing the phosphate from $\text{PEP}$ initially yields the unstable enol tautomer of pyruvate, which spontaneously and rapidly tautomerizes into the far more stable keto form, driving the large negative $\Delta G^{\circ\prime}$.

\item \textbf{(B)} Water is produced only at the Enolase step ($1\text{ H}_2\text{O}$ per triose). Since $2\text{ molecules}$ of 2-phosphoglycerate are dehydrated per glucose, $2\text{ molecules}$ of $\text{H}_2\text{O}$ are produced.

\item \textbf{(D)} Statement (D) is false because glycolysis occurs exclusively in the **cytoplasm (cytosol)** of all living cells, both prokaryotic and eukaryotic—never in the mitochondrial matrix.

\item \textbf{(B)} The phosphate at $\text{C}_1$ of 1,3-bisphosphoglycerate is a high-energy mixed anhydride (acyl phosphate) bond, which is transferred to $\text{ADP}$ by phosphoglycerate kinase. The phosphate at $\text{C}_3$ is a low-energy ester bond retained on 3-phosphoglycerate.

\item \textbf{(B)} Per glucose: $2\text{ ATP}$ are consumed in the preparatory phase and $4\text{ ATP}$ are synthesized in the payoff phase (net gain = $2\text{ ATP}$). For $5\text{ glucose}$ molecules: $5 \times 4 = 20\text{ ATP}$ synthesized, and $5 \times 2 = 10\text{ ATP}$ consumed.

\item \textbf{(B)} Pyruvate ($\alpha$-keto acid) undergoes reversible transamination (via alanine aminotransferase, $\text{ALT}$) with glutamate to form the amino acid **alanine**.

\item \textbf{(A)} At $\text{C}_1$ of $\text{PGAL}$, an aldehyde group ($-\text{CHO}$, oxidation state $+1$) is oxidized to a carboxylic acid derivative—specifically an acyl phosphate ($-\text{COO}-\text{PO}_3^{2-}$, oxidation state $+3$).

\item \textbf{(D)} Statement I is incorrect: Fructose-2,6-bisphosphate is a regulatory signaling molecule synthesized by $\text{PFK-2}$, not a metabolic intermediate of glycolysis (which uses fructose-1,6-bisphosphate). Statement II is correct: it strongly activates $\text{PFK-1}$ in liver cells.

\item \textbf{(B)} Enolase catalyzes the conversion of 2-phosphoglycerate to phosphoenolpyruvate. If Enolase is blocked, its immediate substrate, **2-phosphoglycerate**, will accumulate to the highest levels.

\end{enumerate}

\end{document}